\documentclass[12pt]{article}

\usepackage[utf8]{inputenc}
\usepackage{amsmath, amssymb, amsthm}
\usepackage{geometry}
\geometry{a4paper, margin=2.5cm}

\title{Forma Normală Smith}
\author{}
\date{}

\theoremstyle{definition}
\newtheorem{definition}{Definiție}

\theoremstyle{plain}
\newtheorem{theorem}{Teoremă}

\begin{document}

\maketitle

\begin{definition}
Fie $R$ un domeniu principal (PID). Două matrici $A, B \in M_{m,n}(R)$ se numesc
\emph{aritmetic echivalente} dacă există matrici inversabile


\[
U \in GL_m(R), \qquad V \in GL_n(R)
\]


astfel încât


\[
B = UAV.
\]


\end{definition}

\begin{definition}
Pentru o matrice $A \in M_{m,n}(R)$, notăm prin $\Delta_k$ cel mai mare divizor comun al tuturor minorilor de ordin $k$ ai lui $A$. Acestea se numesc \emph{determinante elementare} ale lui $A$.
\end{definition}

\begin{theorem}[Forma Normală Smith]
Fie $R$ un PID și $A \in M_{m,n}(R)$. Atunci există matrici inversabile


\[
U \in GL_m(R), \qquad V \in GL_n(R)
\]


astfel încât


\[
UAV = D,
\]


unde $D$ este o matrice diagonală de forma


\[
D = \operatorname{diag}(d_1, d_2, \dots, d_r, 0, \dots, 0),
\]


cu proprietățile:
\begin{enumerate}
    \item[(i)] $d_i \ne 0$ pentru $1 \le i \le r$;
    \item[(ii)] $d_1 \mid d_2 \mid \cdots \mid d_r$;
    \item[(iii)] $D$ este unică până la unități din $R$.
\end{enumerate}
Numerele $d_1, \dots, d_r$ se numesc \emph{divizori invarianti} ai lui $A$.
\end{theorem}

\begin{theorem}[Legătura cu determinantele elementare]
În ipotezele teoremei precedente, divizorii invarianti $d_1,\dots,d_r$ satisfac relațiile:


\[
d_1 \sim \Delta_1, \qquad
d_2 \sim \frac{\Delta_2}{\Delta_1}, \qquad
\dots, \qquad
d_r \sim \frac{\Delta_r}{\Delta_{r-1}},
\]


unde $\Delta_k$ este cel mai mare divizor comun al minorilor de ordin $k$ ai lui A, iar $\sim$ înseamnă egalitate până la un element unitate din $R$.
\end{theorem}

\end{document}
